% Authors: Teddy Wachtler
%%%%%%%%%%%%%%%%%%%%%%%%%%%%%%%%%%%%%%%%%%%%%%%

\documentclass{article}
\usepackage[dvips]{graphicx}
\usepackage[a4paper, margin=3cm]{geometry}
\usepackage{amssymb}
\usepackage{amsthm}
\usepackage{amsopn}
\usepackage{amsmath}

\theoremstyle{definition}
\newtheorem*{definition}{Definition}
\newtheorem{theorem}{Theorem}

\newcommand{\header}[3]{
	\begin{center}
	\Large{{#1}} \\
	\smallskip \normalsize {#2} \\
	\smallskip \hfill Name: {#3} \\
	\smallskip \hrule
	\bigskip
	\end{center}}
\newcommand{\problem}[1]{
	\noindent Problem {#1}
	\medskip }
\newcommand{\proposition}{\underline{Proposition: }}
\renewcommand{\proof}{
	\medskip \fbox{Proof}
	\medskip }
\renewcommand{\qed}{\blacksquare}
\newcommand{\todo}{\textcolor{red}{TO-DO.} }

\newcommand{\set}[1]{\{{#1}\}}
\newcommand{\powset}{{\cal P}}
\newcommand{\img}{_{*}}
\newcommand{\pimg}{^{*}}

\newcommand{\vv}{\ensuremath{\vec{v}}}
\newcommand{\vu}{\ensuremath{\vec{u}}}
\newcommand{\vw}{\ensuremath{\vec{w}}}
\newcommand{\vx}{\ensuremath{\vec{x}}}
\newcommand{\vy}{\ensuremath{\vec{y}}}
\newcommand{\vb}{\ensuremath{\vec{b}}}
\newcommand{\vo}{\ensuremath{\vec{0}}}
\newcommand{\va}{\ensuremath{\vec{a}}}
\newcommand{\ve}{\ensuremath{\vec{e}}}

\newcommand{\R}{\mathbb{R}}
\newcommand{\Z}{\mathbb{Z}}
\newcommand{\C}{\mathbb{C}}
\newcommand{\N}{\mathbb{N}}
\newcommand{\Q}{\mathbb{Q}}

%%%%%%%%%%%%%%%%%%%%%%%%%%%%%%%%%%%%%%%%%%%%%%%

\begin{document}
	\header
	{Math 251W: Foundations of Advanced Mathematics}
	{Portfolio Assignment 5: \S 5.1 \& 5.3}
	{Teddy Wachtler}

%%%%%%%%%%%%%%%%%%%%%%%%%%%%%%%%%%%%%%%%%%%%%%%

	\problem{51}

	\proposition
	Let $A$ be a set, and let $R$ be a relation on $A$.

	\begin{enumerate}
		\item
		Suppose that $R$ is reflexive.
		Prove that $\bigcup_{x \in A} [x] = A$.

		\proof

		Let $A$ be an arbitrary set.
		Let $R$ be an arbitrary relation on $A$ so that $R$ is reflexive.

		Let $a_0$ be an arbitrary element of $\bigcup_{x \in A} R[x]$.
		By the definition of indexed unions, we see that there exists an element $b \in A$ such that $a_0 \in R[b]$.
		By the definition of the equivalence class, we see that $R[b] \subseteq A$.
		By the definition of subsets, since $a_0 \in R[b]$, we see that $a_0 \in A$.
		Therefore, by the definition of subsets, since $a_0 \in \bigcup_{x \in A} R[x]$ implies $a_0 \in A$, we see that $\bigcup_{x \in A} R[x] \subseteq A$.

		Let $a_1$ be an arbitrary element of $A$.
		By the definition of reflexive relations, we see that $(a_1, a_1) \in \bar R$.
		By the definition of the equivalence class, we see that $a_1 \in R[a_1]$.
		By the definition of indexed unions, since there exists $a_1 \in A$ so that $a_1 \in R[a_1]$, we see that $a_1 \in \bigcup_{x \in A} R[x]$.
		Therefore, by the definition of subsets, since $a_1 \in A$ implies $a_1 \in \bigcup_{x \in A} R[x]$, we see that $A \subseteq \bigcup_{x \in A} R[x]$.

		By the definition of set equality, since $\bigcup_{x \in A} R[x] \subseteq A$ and $A \subseteq \bigcup_{x \in A} R[x]$, we see that $\bigcup_{x \in A} R[x] = A$.
		Therefore, we see that for all sets $A$ and for all relations $R$ on $A$ so that $R$ is reflexive, it is the case that $\bigcup_{x \in A} R[x] = A$.
		$\qed$.

		\smallskip

		\item
		Suppose that $R$ is symmetric.
		Prove that $x \in [y]$ if and only if $y\in [x]$, for all $x, y \in A$.

		\proof

		Let $A$ be an arbitrary set.
		Let $R$ be an arbitrary relation on $A$ so that $R$ is symmetric.
		Let $x$ and $y$ be arbitrary elements of $A$.
		Suppose $x \in R[y]$.
		By the definition of the equivalence class, we see that $(x, y) \in \bar R$.
		By the definition of symmetric relations, since $(x, y) \in \bar R$, we see that $(y, x) \in \bar R$.
		By the definition of the equivalence class, we see that $y \in R[x]$.
		Therefore, we see that if $x \in R[y]$, then $y \in R[x]$.
		By a similar argument, exchanging $x$ for $y$, we see that if $y \in R[x]$, then $x \in R[y]$.
		Since if $x \in R[y]$, then $y \in R[x]$, and if $y \in R[x]$, then $x \in R[y]$, we see that $x \in R[y]$ if and only if $y \in R[x]$.
		Therefore, we see that for all sets $A$, for all relations $R$ on $A$ so that $R$ is symmetric, and for all elements $x$ and $y$ in $A$, it is the case that $x \in R[y]$ if and only if $y\in R[x]$.
		$\qed$.

		\smallskip

		\item
		Suppose that $R$ is transitive.
		Prove that if $x R y$, then $[y] \subseteq [x]$, for all $x, y \in A$.

		\proof

		Let $A$ be an arbitrary set.
		Let $R$ be an arbitrary relation on $A$ so that $R$ is transitive.
		Let $x$ and $y$ be arbitrary elements of $A$.
		Suppose $(x, y) \in \bar R$.
		Let $a$ be an arbitrary element of $R[y]$.
		By the definition of the equivalence class, we see that $(y, a) \in \bar R$.
		By the definition of transitive relations, since $(x, y) \in \bar R$ and  $(y, a) \in \bar R$, we see that $(x, a) \in \bar R$.
		By the definition of the equivalence class, we see that $a \in R[x]$.
		By the definition of subsets, since $a \in R[y]$ implies $a \in R[x]$, we see that $R[y] \subseteq R[x]$.
		Therefore, we see that for all sets $A$, for all relations $R$ on $A$ so that $R$ is transitive, and for all elements $x$ and $y$ in $A$, it is the case that if $(x, y) \in \bar R$, then $R[y] \subseteq R[x]$.
		$\qed$.

		\smallskip

	\end{enumerate}

%%%%%%%%%%%%%%%%%%%%%%%%%%%%%%%%%%%%%%%%%%%%%%%

\end{document}